% !TeX root = ../main.tex

\chapter{总结与展望}
\section{研究总结}
在信息化教育时代的历史推进下，国家教育政策要求将信息科技深度融入数学教学的过程。如何适当地选择辅助教学的实验工具，利用信息技术软件进行对数学实验教学资源的设计与开发，学生完成从“听数学”到“做数学”的学习状态转变，是本研究的主要目标。对此，本研究的核心是基于GeoGebra软件和希沃白板软件，开发部分高中数学内容模块的实验教学资源，开展了以下主要研究工作：
\begin{enumerate}[(1)]
  \item 研读近年来国家颁布的教育指导文件，及高中数学课程标准和义务教育数学课程标准，明晰研究的社会背景。查阅相关文献，对国内外数学实验教学的研究进展进行梳理。整理发现，多数研究是核心探索实验在数学教学中的应用，少有实验教学资源开发类研究。据此，明确本研究的目的，即开发高中数学实验教学资源，并配置相关片段实验讲义。
  \item 梳理研究脉络、查阅相关文献，阐述本研究的目的与意义、研究内容和研究方法，对涉及的概念进行明细界定，同时介绍教育与心理理论支撑。对研究进展的必要性进行开发分析，确保本研究的合理性与被需性。
  \item 基于GeoGebra软件和希沃白板软件进行高中数学实验教学资源的设计与开发，分步骤详细说明资源开发过程，以便于读者开展自主复制设计与二次开发。此外，设计配套片段实验讲义，作为教学资源元件供教学设计嵌入式使用，可在此基础上适当改动以适配不同教学模式。
\end{enumerate}

本研究的创新价值体现在多重维度突破：技术层面，通过GeoGebra脚本引擎开发一键式操作界面，降低动态数学实验的技术门槛，使教师无需复杂编程即可生成函数图象连续变换动画；内容层面，针对人教A版必修教材“三角函数的图象与性质”、“平面向量基本定理”等章节，设计师生交互式实验演示实验，将操作过程动态可视化，不再局限于口头讲述与空间无支撑猜想，实现信息化赋能课堂教学；资源层面：数字化实验教学资源可以长期保存、动态更新，便于追踪教学效果，总结反思进而改进，沉淀成系统性知识库；发展层面，提出自主创作实验教学资源共享化，即建立私有资源共享网页，打破地域限制，让广大师生获取优质资源的同时，减少教师重复备课的负担，促进教师间教学经验交流与创新，形成“共建—优化—再应用”的良性循环。

本研究存在以下主要不足：
\begin{enumerate}[(1)]
  \item 对GeoGebra和希沃白板的高级功能及深度应用技巧有待深化掌握。例如，GeoGebra中复杂脚本语言的运用以及希沃白板在跨学科融合场景下的拓展功能，尚未能充分挖掘和熟练驾驭。
  \item 受限于时间与精力，目前仅利用GeoGebra和希沃白板开发了代数、几何部分的部分典型实验。像概率统计、数学建模等领域的实验开发相对匮乏，开发案例不够全面，难以覆盖高中数学的完整知识体系，后续还有大量的实验开发工作有待推进。
  \item 论文在教育实证方面有所欠缺。因个人教学实践时间紧凑和实践环境的局限，仅对已开发的数学实验案例完成了片段实验讲义，未将其切实应用于真实课堂教学实践中，深入探究这些实验在真实课堂环境下对学生学习效果、思维培养以及数学核心素养提升所产生的实际影响，这无疑是后续研究工作需重点攻坚和突破的关键方向。
\end{enumerate}


\section{研究展望}

信息技术的迅猛发展为高中数学教学开辟了全新的创新路径。本研究聚焦于如何淋漓尽致地彰显自主开发特性，充分发挥 GeoGebra 和希沃白板的独特优势，深入探索二者与高中数学教学的深度融合。旨在降低数学实验操作门槛，丰富教师教学与学生学习数学知识的数字化手段。​

在研究内容维度，针对现有的基于 GeoGebra 和希沃白板开发的数学实验教学资源，有必要深入挖掘其潜在功能，将其广泛拓展至更多样化的数学实验设计与探究中。例如，将三角函数图象变换实验延伸至函数性质的综合探究，或是把圆锥曲线形成实验应用于解析几何不同课型的实际课题里，全方位展现其更为广阔的教育价值。在研究工具层面，GeoGebra 以其强大的动态数学功能为数学实验提供了灵活多变的模型构建能力，希沃白板则凭借出色的多媒体展示与互动特性为课堂教学增添了丰富的交互体验。

未来研究需要着力探索如何在不同的数学课堂情境中，精准且恰当地选择信息化技术工具、自主且完善地开发实验式教学资源，使其能够切实有效地促进学生数学核心素养的全面提升，这仍需进一步深入研究与持续探索。